\subsection{Video Editing}
\label{subsec:p3_s8_3_5}

Video editing takes an \emph{existing} clip and modifies it under user instruction, rather than generating one from scratch (§\ref{subsec:p3_s8_3_1}) or steering that generation with controls (§\ref{subsec:p3_s8_3_2}--§\ref{subsec:p3_s8_3_4}). It is the temporal lift of the image editing of §\ref{subsec:p3_s8_1_3}, and its defining demand is selective preservation: the operator must change only the targeted content while holding the source's geometry, layout, and motion fixed across every frame, so the central technical problem is to balance editability against spatio-temporal consistency without artifacts or unintended drift.

\subsubsection{Task Definition}
\label{subsubsec:p3_s8_3_5_task}
Given a source clip encoded as $\mathbf{z}^{\text{in},1:T}=E(\mathbf{v}^{\text{in}})$ and an edit instruction $\mathbf{c}$, an editing operator
\begin{equation}
\hat{\mathbf{z}}^{1:T}=E_\theta^{\text{vid}}\!\big(\mathbf{z}^{\text{in},1:T},\,\phi(\mathbf{x}),\,\kappa(\mathbf{c})\big)
\label{eq:p3_s8_3_5_objective}
\end{equation}
returns an edited clip that applies the requested change while preserving unedited content and the source motion, recovering the image editing operator of Eq.~\eqref{eq:t2i_edit_objective} exactly at $T=1$. Eq.~\eqref{eq:p3_s8_3_5_objective} encodes two competing objectives: \emph{editability} (the output must realize the instruction $\mathbf{c}$) and \emph{preservation} (the unedited regions and the source dynamics must survive untouched).

\subsubsection{Methodology}
\label{subsubsec:p3_s8_3_5_methodology}
The five paradigms are overlapping choices of training strategy and preservation mechanism: training-free editing keeps the weights fixed, one-shot tuning updates selected parameters on the source clip, trained instruction editors learn across editing data, attention-based editing reuses or modifies source attention, and structural guidance conditions the edit on source-derived depth, edges, or poses. Many of them share a diffusion recipe that inverts the source clip and re-denoises it under the instruction; writing $\mathrm{Inv}$ for DDIM inversion of the pretrained or adapted backbone $\epsilon_\theta^{0}$ of §\ref{subsec:p3_s8_3_1},
\begin{equation}
\mathbf{z}_S^{1:T}=\mathrm{Inv}\big(\mathbf{z}^{\text{in},1:T};\,\phi(\mathbf{x}^{\text{src}})\big),\qquad
\hat{\mathbf{z}}^{1:T}=\mathrm{Denoise}\big(\mathbf{z}_S^{1:T};\,\phi(\mathbf{x}),\,\kappa(\mathbf{c})\big),
\label{eq:p3_s8_3_5_invert}
\end{equation}
where $S$ denotes the terminal diffusion step, distinct from the clip length $T$. Inversion supplies a source-correlated initialization; how much of the source survives then depends on the sampling, feature reuse, and guidance applied during re-denoising. Figure~\ref{fig:p3_s8_3_5_mechanisms} schematizes the five paradigms, with panels (a)--(e) keyed to the paragraphs below.

\begin{figure}[!th]
\centering
\includegraphics[width=\linewidth]{figures/part3/section8/five_video_editing_paradigms.pdf}
\caption{Five overlapping video-editing paradigms, (a)--(e), applied to one task (top): restyling a driving clip while preserving its layout and motion. The dashed link in (c) marks a paired training example; scene sequences are illustrative.}
\label{fig:p3_s8_3_5_mechanisms}
\end{figure}

\noindent\textbf{Training-Free (Zero-Shot) Methods.} The lightest realization of Eq.~\eqref{eq:p3_s8_3_5_invert} (Figure~\ref{fig:p3_s8_3_5_mechanisms}a) edits a video without any fine-tuning, analogous to zero-shot image editing~\cite{ju2023directinversionboostingdiffusionbased}: DDIM inversion encodes the source into latent codes and noise maps, and the denoising process is rerun on them under the edited prompt. Text2Video-Zero~\cite{Khachatryan_2023_ICCV} and Pix2Video~\cite{ceylan2023pix2videovideoeditingusing} established the category, and vid2vid-zero~\cite{wang2024zeroshotvideoeditingusing}, I2VEdit~\cite{ouyang2024i2veditfirstframeguidedvideoediting}, RAVE~\cite{kara2024rave}, and FastVideoEdit~\cite{zhang2024fastvideoeditleveragingconsistencymodels} add latent-code warping and feature injection to improve temporal consistency without fine-tuning.

\noindent\textbf{One-Shot Tuning Methods.} Trading a little training for sharper fidelity (Figure~\ref{fig:p3_s8_3_5_mechanisms}b), Tune-A-Video~\cite{wu2023tune} fine-tunes a limited subset of a pretrained text-to-image diffusion model's parameters on the frames of a single source video, so that the model learns its content and motion and can regenerate it under a new prompt. Dreamix~\cite{molad2023dreamixvideodiffusionmodels} builds on one-shot fine-tuning to improve both editability and fidelity to the original motion.

\noindent\textbf{Training-Based Methods.} The limited control and quality of inference-time editing have motivated editors that learn a reusable operator from editing data (Figure~\ref{fig:p3_s8_3_5_mechanisms}c), particularly for instruction-based editing. InsV2V~\cite{cheng2023consistentvideotovideotransferusing} extends the InstructPix2Pix~\cite{brooks2023instructpix2pixlearningfollowimage} framework directly to video; GenProp~\cite{Liu_2025_CVPR} propagates an edit from a user-modified first frame to the rest of the sequence; and VideoPainter~\cite{videopainter} uses a video inpainting model as an intermediary. UNIC~\cite{ye2025unicunifiedincontextvideo} handles many editing tasks in one model by concatenating the conditions into a sequence with task-aware positional embeddings, and EditVerse~\cite{ju2025editverseunifyingimagevideo} unifies video generation and editing by representing text, image, and video as one token sequence processed with self-attention, which permits in-context learning and flexible inputs.

\noindent\textbf{Attention-Based Editing.} To protect preservation within the inference-time recipe (Figure~\ref{fig:p3_s8_3_5_mechanisms}d), many techniques manipulate the cross-attention maps directly: following Prompt-to-Prompt in images, they blend the attention maps of the original and edited prompts so that the inverted source structure of Eq.~\eqref{eq:p3_s8_3_5_invert} survives the re-denoising. FateZero~\cite{qi2023fatezerofusingattentionszeroshot} and Video-P2P~\cite{liu2023videop2p} were early examples, and Edit-A-Video~\cite{shin2023editavideosinglevideoediting}, TokenFlow~\cite{tokenflow2023}, UniEdit~\cite{bai2024uniedit}, and AnyV2V~\cite{ku2024anyv2vtuningfreeframeworkvideotovideo} improve temporal consistency at the token or feature level, enabling localized edits such as changing an object's material or color while the scene structure is kept.

\noindent\textbf{Editing with Structural Guidance.} At the explicitly guided end (Figure~\ref{fig:p3_s8_3_5_mechanisms}e), editing is conditioned on spatial signals extracted from the source video, the conditioning machinery of §\ref{subsec:p3_s8_3_2} repurposed for preservation rather than generation. Rerender A Video~\cite{yang2023rerender} feeds the source's depth or edge maps to a ControlNet-like module during editing so that geometry and layout stay intact, and MagicEdit~\cite{liew2023magicedit} accepts a variety of spatial inputs (\eg, pose, sketches, segmentation maps) for high structural fidelity, which suits style transfer and object replacement, where preserving structure is what prevents distortion and keeps motion coherent.

The paradigms trade cost against control: training-free editing avoids per-clip weight updates but may still need iterative inference, whereas a trained instruction editor is expensive to build but generalizes across edits.

\subsubsection{Benchmarks}
\label{subsubsec:p3_s8_3_5_benchmarks}
Video editing has no single canonical suite, and evaluation has moved from shared source-video collections toward purpose-built, fine-grained editing benchmarks. Early evaluations report on DAVIS~\cite{pont20172017}, the de facto source set for object-centric edits, and on the LOVEU text-guided video-editing (TGVE) benchmark~\cite{wu2023loveutgve}, which supplies prompt--clip pairs with human ratings; VPBench, released with VideoPainter~\cite{videopainter}, adds masked clips with dense captions for inpainting-style edits. The most recent suites target instruction-guided editing and pair richer data with model-based scoring. FiVE~\cite{li2025five} contributes $74$ real clips (drawn from DAVIS) and $26$ generated clips across six fine-grained editing types, with $420$ object-level prompt pairs and masks, evaluated under a $15$-metric protocol spanning background preservation, text--video similarity, temporal consistency, video quality, and runtime. IVEBench~\cite{chen2026ivebench} scales this to $600$ source videos spanning seven semantic dimensions and clip lengths from $32$ to $1{,}024$ frames, across eight task categories ($35$ subcategories) whose prompts are generated by LLMs and refined by experts, and adopts a three-dimensional protocol (video quality, instruction compliance, and video fidelity) that fuses traditional metrics with multimodal-LLM-based assessment.

\paragraph{Editing-fidelity Metrics.}
Editing is judged on three axes that mirror the editability--preservation tension of Eq.~\eqref{eq:p3_s8_3_5_objective}. The classic metrics are, for \emph{edit adherence}, the CLIP similarity between each edited frame and the target prompt; for \emph{temporal consistency}, frame-to-frame CLIP similarity together with a warping error that flow-aligns adjacent frames and reports the photometric residual; and for \emph{structure and background preservation}, PSNR / SSIM / LPIPS on the unedited regions against the source, atop the generic FVD~\cite{unterthiner2018fvd} quality score. These scalars correlate only loosely with human judgment on fine-grained edits (an object recolored but mislocalized can still score well), so the latest benchmarks add vision-language judges: FiVE-Acc~\cite{li2025five} uses a VLM to test whether the intended fine-grained edit succeeded, and IVEBench~\cite{chen2026ivebench} scores instruction compliance and fidelity with multimodal LLMs alongside the traditional metrics. Results are reported as an edit-versus-preservation pair, since ignoring the instruction trivially maximizes preservation.

\subsubsection{Challenges and Trends}
\label{subsubsec:p3_s8_3_5_challenges}
Three problems are specific to editing. Pushing the edit harder erodes the unedited regions and the source motion, while protecting them too strongly suppresses the requested change; this editability--preservation trade-off of Eq.~\eqref{eq:p3_s8_3_5_objective} is what every paradigm above negotiates differently. The invert-and-re-denoise recipe of Eq.~\eqref{eq:p3_s8_3_5_invert} is imperfect, and its reconstruction error compounds over long clips, so edits that are clean over a few seconds degrade over minute-scale video. And training-based editors follow only the edit types seen in their training pairs, so a flexible instruction editor that handles arbitrary edits remains open; the unified generate-and-edit models UNIC~\cite{ye2025unicunifiedincontextvideo} and EditVerse~\cite{ju2025editverseunifyingimagevideo}, which treat the source clip as one stream among many in shared attention, are the current attempt.
